%% ma: L10-B15-0081
%% lop: 10
%% chuong: 6
%% bai: 15
%% dang: 10.15.1
%% loai: DS
%% muc: [NB, TH, TH, VD]
%% dap_an: "ĐSĐĐ"
%% nguon: "Ảnh giáo viên gửi 25/11/2026: Bài tập trắc nghiệm – Toán 10 Bài 15 Hàm số (3 phần), Phần II Câu 3"
%% kien_thuc: "Lớp 10 Bài 15 (tập xác định của hàm số); bất phương trình bậc nhất, giao các khoảng (Lớp 10 Bài 2)"
%% trang_thai: chinh-thuc
\begin{ex}
Cho hàm số $f(x)=\sqrt{a+2x}+\sqrt{a+3-x}$ (với $a$ là tham số). Xét tính đúng sai của các khẳng định sau:
\choiceTF
{\True Hàm số xác định khi và chỉ khi $\begin{cases}a+2x\ge0\\ a+3-x\ge0\end{cases}$.}
{Với $a=1$, hàm số xác định khi và chỉ khi $\begin{cases}1+2x\ge0\\ 2-x\ge0\end{cases}$.}
{\True Với $a=2$, tập xác định của hàm số là $D=[-1;5]$.}
{\True Với $a=3$, tập xác định của hàm số là $D$. Khi đó tập hợp $D\cap[1;8]$ có $6$ giá trị nguyên.}
\loigiai{a) Đúng. b) Với $a=1$ điều kiện thứ hai là $4-x\ge0$: Sai. c) $2+2x\ge0$, $5-x\ge0\Leftrightarrow-1\le x\le5$: Đúng. d) $D=\left[-\tfrac32;6\right]$, $D\cap[1;8]=[1;6]$ có $6$ số nguyên: Đúng.}
\end{ex}
